\documentclass[12pt,a4paper]{article}

% ---------------------------------------------------------------------------
% Personalized Research Paper Recommendation System -- Project Report
% Style mirrors the provided sample report (article class, Computer Modern,
% numbered sections, simple ruled tables).
% ---------------------------------------------------------------------------

\usepackage[utf8]{inputenc}
\usepackage[T1]{fontenc}
\usepackage[english]{babel}
\usepackage{lmodern}
\usepackage[margin=1in]{geometry}
\usepackage{amsmath,amssymb}
\usepackage{graphicx}
\usepackage{booktabs}
\usepackage{array}
\usepackage{enumitem}
\usepackage{caption}
\usepackage{float}
\usepackage{xcolor}
\usepackage[hidelinks]{hyperref}

\captionsetup{font=small,labelfont=bf}
\setlength{\parskip}{0.4em}
\setlength{\parindent}{0pt}
\renewcommand{\arraystretch}{1.25}

% Safe figure inclusion: if the image file is missing, draw a labelled
% placeholder box so the document always compiles inside or outside the repo.
\newcommand{\safefigure}[3]{%
  % #1 = path, #2 = caption, #3 = label
  \begin{figure}[H]
    \centering
    \IfFileExists{#1}{%
      \includegraphics[width=0.85\linewidth]{#1}%
    }{%
      \fbox{\parbox[c][4.5cm][c]{0.85\linewidth}{\centering\itshape
        Figure placeholder --- expected image:\\\texttt{\detokenize{#1}}\\[2pt]
        (rendered when compiled inside the project repository)}}%
    }
    \caption{#2}\label{#3}
  \end{figure}
}

\title{\bfseries Personalized Research Paper Recommendation System\\[6pt]
\large with Temporal Scoring, Diversity-aware Ranking,\\
Learned Ranking, and Fine-tuned Semantic Embeddings\\[10pt]
\normalsize Machine Learning Project Report --- IT3190E}
\author{}
\date{June 22, 2026}

\begin{document}
\maketitle
\thispagestyle{empty}

% ===========================================================================
\section{Introduction}

\subsection{General Overview}
This project builds a personalized, content-based recommendation system for
research papers that combines semantic NLP embeddings with temporal awareness
and diversity-conscious ranking. The system models a user's research interests
from their reading history and then returns recommendations that are
simultaneously relevant, recent, and non-redundant.

The central idea is that a good recommendation balances three factors: the
\emph{semantic relevance} of a paper to the user's interests, the
\emph{recency} of its publication, and the \emph{topical diversity} of the
returned list. The system addresses all three through a hybrid scoring
pipeline built on Sentence-BERT (SBERT) embeddings, exponential temporal
scoring, and the Maximal Marginal Relevance (MMR) algorithm.

Beyond a baseline retrieval pipeline, the project introduces two learned
components that lift the system beyond plain similarity search:
\begin{itemize}[leftmargin=1.4em]
  \item \textbf{Learning-to-Rank (LTR):} a Gradient Boosting model that learns
  how to combine ranking signals instead of hard-coding weighting parameters.
  \item \textbf{Fine-tuned SBERT:} domain-adapted embeddings produced by
  contrastive learning on arXiv papers, so that papers with similar content
  are closer in embedding space.
\end{itemize}

\subsection{Problem Statement}
Given a corpus of research papers and a user described only by an implicit
reading history, how can the system rank papers so that the top results are
relevant to the user's current interests, favour recently published work, and
avoid returning near-duplicate items? No single fixed scoring rule is optimal:
pure semantic similarity ignores recency and tends to produce redundant lists,
while naive recency ignores meaning. The pipeline therefore composes several
complementary signals and a learned ranker.

\subsection{Project Objectives}
The main objectives, following the project proposal, are to:
\begin{enumerate}[leftmargin=1.6em]
  \item Build a content-based recommendation system with semantic understanding
  of paper abstracts.
  \item Model user interests from reading history as a temporally-weighted
  embedding profile.
  \item Implement and compare TF-IDF, pretrained SBERT, and fine-tuned SBERT
  representations.
  \item Incorporate publication recency into the recommendation score via
  exponential decay.
  \item Improve recommendation diversity with Maximal Marginal Relevance.
  \item Train a Learning-to-Rank model that learns ranking signals from
  simulated interactions.
  \item Fine-tune SBERT through contrastive learning on arXiv papers.
  \item Handle cold-start users through category onboarding.
  \item Evaluate using Precision@K, NDCG@K, Intra-list Diversity, and Semantic
  Coherence.
\end{enumerate}

The proposed pipeline is summarised below:
\begin{center}
\fbox{\parbox{0.9\linewidth}{\ttfamily\small
Preprocessing\\
\hspace*{1em}$\rightarrow$ TF-IDF / SBERT / Fine-tuned SBERT\\
\hspace*{1em}$\rightarrow$ Temporal-weighted user profile\\
\hspace*{1em}$\rightarrow$ Hybrid scoring or Learning-to-Rank\\
\hspace*{1em}$\rightarrow$ MMR reranking\\
\hspace*{1em}$\rightarrow$ Explainable Top-K recommendations}}
\end{center}

% ===========================================================================
\section{Dataset}

The proposal targets a corpus of 7{,}701 AI/ML arXiv papers spanning
September~2025 to April~2026. The current repository extends this into two
splits, summarised in Table~\ref{tab:dataset}.

\begin{table}[H]
\centering
\caption{Dataset splits used in the project.}\label{tab:dataset}
\begin{tabular}{lrl}
\toprule
\textbf{Split} & \textbf{Papers} & \textbf{Purpose}\\
\midrule
\texttt{arxiv\_dataset/train.csv} & 136{,}238 & Embeddings, user profiles, fine-tuning, LTR\\
\texttt{arxiv\_dataset/test.csv}  & 7{,}701   & Representation evaluation and candidate pool\\
\bottomrule
\end{tabular}
\end{table}

The processed CSV files retain the following fields:
\begin{center}\ttfamily\small
paper\_id, title, primary\_category, authors, first\_author,\\
abstract, days\_since\_submission, recency\_score
\end{center}

Relative to the proposal, the current data does not directly preserve the
fields \texttt{all\_categories}, \texttt{submitted\_date},
\texttt{update\_lag\_days}, \texttt{is\_updated}, and
\texttt{is\_large\_collaboration}. This matters for the LTR component: the
proposal specifies five features
(\texttt{cosine\_sim}, \texttt{recency\_score}, \texttt{update\_lag\_days},
\texttt{is\_updated}, \texttt{category\_match}), whereas the current
implementation uses three
(\texttt{cosine\_sim}, \texttt{recency\_score}, \texttt{category\_match}).

% ===========================================================================
\section{Data Preprocessing}

Preprocessing is implemented in \texttt{preprocess.py} /
\texttt{preprocess.ipynb}, and the fitted text model is stored at
\texttt{arxiv\_dataset/text\_preprocessor.pkl}. Abstracts are cleaned of
LaTeX markup, lowercased, and tokenized. The representation artifacts produced
are listed in Table~\ref{tab:repr-artifacts}.

\begin{table}[H]
\centering
\caption{Representation artifacts (matrix shapes).}\label{tab:repr-artifacts}
\begin{tabular}{lrr}
\toprule
\textbf{Representation} & \textbf{Train} & \textbf{Test}\\
\midrule
TF-IDF          & $(136238,\,10000)$ & $7{,}701$ rows\\
SBERT pretrained & $(136238,\,384)$   & $(7701,\,384)$\\
SBERT fine-tuned & $(136238,\,384)$   & $(7701,\,384)$\\
\bottomrule
\end{tabular}
\end{table}

The recency score uses an exponential decay with a configurable half-life:
\begin{equation}
\text{recency\_score} = \exp\!\left(-\ln 2 \cdot \frac{\text{age\_in\_days}}{\text{half\_life}}\right).
\end{equation}

\safefigure{temporal_weight_decay.png}{Temporal weight decay for several
half-life settings.}{fig:decay}

\begin{table}[H]
\centering
\caption{Decay weight as a function of half-life and elapsed time.}\label{tab:decay}
\begin{tabular}{rrrrr}
\toprule
\textbf{Half-life} & \textbf{30 days} & \textbf{60 days} & \textbf{90 days} & \textbf{180 days}\\
\midrule
30 days & 0.5000 & 0.2500 & 0.1250 & 0.0156\\
60 days & 0.7071 & 0.5000 & 0.3536 & 0.1250\\
90 days & 0.7937 & 0.6300 & 0.5000 & 0.2500\\
\bottomrule
\end{tabular}
\end{table}

A smaller half-life makes older behaviour lose influence faster. A 90-day
half-life retains long-term interests better and is used as the primary
evaluation configuration.

% ===========================================================================
\section{Text Representation}

\subsection{TF-IDF}
TF-IDF is a 10{,}000-dimensional bag-of-words baseline. It is simple,
keyword-explainable, and tends to produce highly diverse lists. Its main
weakness is that it does not capture semantic relationships between different
words that share meaning.

\subsection{Pretrained SBERT}
The \texttt{all-MiniLM-L6-v2} model produces dense 384-dimensional vectors.
Pretrained SBERT strongly improves semantic similarity over TF-IDF without any
additional training.

\subsection{Fine-tuned SBERT}
The repository contains two checkpoints:
\begin{center}\ttfamily\small
finetune\_BERT/model\_checkpoint/finetuned-arxiv-sbert\_v1\\
finetune\_BERT/model\_checkpoint/finetuned-arxiv-sbert\_v2
\end{center}

The proposal describes fine-tuning with \texttt{CosineSimilarityLoss} on
positive/negative pairs. The current implementation instead uses
\texttt{TripletLoss} with the following configuration:
\begin{itemize}[leftmargin=1.4em]
  \item base model: \texttt{all-MiniLM-L6-v2};
  \item cosine distance;
  \item triplet margin: \texttt{0.5};
  \item batch size: \texttt{16};
  \item epochs: \texttt{3}.
\end{itemize}
This is a deliberate change from the proposal. Triplet loss still satisfies the
contrastive-learning objective, but the report records it explicitly to avoid
misdescribing the method that was actually run. The triplet-generation script
supports several positive/negative strategies, but the exact triplet input file
it references is not present in the repository; checkpoints and output
embeddings exist, so the pipeline is not yet fully reproducible from the stored
files alone.

% ===========================================================================
\section{Temporal-Weighted User Profile}

A user vector is computed from their reading history with decay weighting:
\begin{align}
\text{weight}(p) &= \exp\!\left(-\ln 2 \cdot \frac{\text{days\_since\_read}(p)}{\text{half\_life}}\right),\\
\text{user\_vector} &= \text{weighted\_mean}\big(\text{paper\_embeddings},\ \text{weights}\big).
\end{align}

Profiles were generated for 1{,}000 synthetic users at half-life settings of
30, 60, and 90 days, for both TF-IDF and SBERT.

\safefigure{user_profile_statistics.png}{User profile vector statistics.}{fig:profile}

\begin{table}[H]
\centering
\caption{TF-IDF user-profile vector norms.}\label{tab:tfidf-norm}
\begin{tabular}{rrrrrr}
\toprule
\textbf{Half-life} & \textbf{Min} & \textbf{Max} & \textbf{Mean} & \textbf{Median} & \textbf{Std}\\
\midrule
30 & 0.2087 & 0.9442 & 0.3403 & 0.2973 & 0.1323\\
60 & 0.2034 & 0.8453 & 0.3330 & 0.2908 & 0.1267\\
90 & 0.2025 & 0.8038 & 0.3317 & 0.2901 & 0.1253\\
\bottomrule
\end{tabular}
\end{table}

\begin{table}[H]
\centering
\caption{SBERT user-profile vector norms.}\label{tab:sbert-norm}
\begin{tabular}{rrrrrr}
\toprule
\textbf{Half-life} & \textbf{Min} & \textbf{Max} & \textbf{Mean} & \textbf{Median} & \textbf{Std}\\
\midrule
30 & 0.1903 & 0.9749 & 0.4105 & 0.3839 & 0.1300\\
60 & 0.2052 & 0.9748 & 0.4233 & 0.4054 & 0.1234\\
90 & 0.2101 & 0.9748 & 0.4289 & 0.4138 & 0.1217\\
\bottomrule
\end{tabular}
\end{table}

As the half-life grows, the average TF-IDF profile norm decreases slightly
while the average SBERT profile norm increases. Norm, however, does not
directly measure recommendation quality. The proposal calls for comparing the
mean profile against the decay-weighted profile using ranking metrics; this
ablation is not yet quantified.

% ===========================================================================
\section{Cold-Start Handling}

Cold-start is implemented in \texttt{cold\_start.py} and integrated into the
Streamlit interface. The procedure is:
\begin{enumerate}[leftmargin=1.6em]
  \item the user selects one to three categories;
  \item for each category, the system takes up to 20 papers with the highest
  \texttt{recency\_score};
  \item the cold-start vector is the mean embedding of the selected papers;
  \item this vector is passed through LTR and MMR exactly like a user with
  history.
\end{enumerate}

\begin{center}\ttfamily\small
cold\_start\_vector =\\
\hspace*{1em}mean(embeddings of top-20 recent papers per selected category)
\end{center}

Unit tests covering selection of the 20 most recent papers, categories with
fewer than 20 papers, duplicate categories, and invalid input all pass.

% ===========================================================================
\section{Hybrid Scoring}

The hybrid score linearly combines semantic similarity and recency:
\begin{equation}
\text{final\_score} = \alpha \cdot \text{cosine\_similarity}
+ (1-\alpha)\cdot \text{recency\_score}.
\end{equation}
The primary configuration uses $\alpha = 0.7$, half-life $=90$, and Top-$K=10$.
The repository stores results for 1{,}000 users at half-life 30, 60, and 90.

\safefigure{recommendations/hybrid_scoring_analysis.png}{TF-IDF hybrid scoring analysis.}{fig:hybrid-tfidf}
\safefigure{recommendations_SBERT/hybrid_scoring_analysis.png}{SBERT hybrid scoring analysis.}{fig:hybrid-sbert}

\begin{table}[H]
\centering
\caption{Hybrid scoring means computed directly from the stored JSON.}\label{tab:hybrid}
\begin{tabular}{lrrrr}
\toprule
\textbf{Representation} & \textbf{Half-life} & \textbf{Mean sim.} & \textbf{Mean recency} & \textbf{Mean final}\\
\midrule
TF-IDF & 30 & 0.2532 & 0.0506 & 0.1924\\
TF-IDF & 60 & 0.2500 & 0.0505 & 0.1902\\
TF-IDF & 90 & 0.2492 & 0.0505 & 0.1896\\
SBERT  & 30 & 0.7123 & 0.0505 & 0.5138\\
SBERT  & 60 & 0.7107 & 0.0505 & 0.5126\\
SBERT  & 90 & 0.7103 & 0.0505 & 0.5124\\
\bottomrule
\end{tabular}
\end{table}

\begin{table}[H]
\centering
\caption{Actual contribution share at half-life 90.}\label{tab:contrib}
\begin{tabular}{lrr}
\toprule
\textbf{Representation} & \textbf{Semantic} & \textbf{Recency}\\
\midrule
TF-IDF & 92.00\% & 8.00\%\\
SBERT  & 97.04\% & 2.96\%\\
\bottomrule
\end{tabular}
\end{table}

Although the configured weights are 70\% semantic and 30\% recency, the two
components are not on the same scale: recommendation \texttt{recency\_score}
values cluster around $0.05$, so the real recency contribution is only about
3--8\%. To make recency influence match the configured ratio, the two features
should be normalised or calibrated before being combined. The two hybrid PNG
figures appear to have been generated from runs with a different configuration
than the current folder names, so the table above (sourced from JSON) is
treated as the authoritative quantitative source.

% ===========================================================================
\section{Learning-to-Rank}

The LTR component is a \texttt{GradientBoostingClassifier}. Training data are
generated from synthetic users via leave-one-out simulation:
\begin{itemize}[leftmargin=1.4em]
  \item positive: the hidden target paper;
  \item negative: sampled non-relevant papers;
  \item labels: 1 for positive, 0 for negative.
\end{itemize}
Two models are stored, \texttt{ltr\_tfidf.pkl} and \texttt{ltr\_sbert.pkl}, and
LTR is integrated into the application to score the full candidate set before
MMR reranking.

\begin{table}[H]
\centering
\caption{Differences between proposal and current LTR implementation.}\label{tab:ltr-diff}
\begin{tabular}{ll}
\toprule
\textbf{Proposal} & \textbf{Current implementation}\\
\midrule
5 features & 3 features\\
Includes \texttt{update\_lag\_days} & Not included\\
Includes \texttt{is\_updated} & Not included\\
LTR vs cosine baseline by NDCG & No benchmark table yet\\
\bottomrule
\end{tabular}
\end{table}

The model artifacts exist and are used in the UI, but there is not yet a
quantitative result demonstrating that LTR outperforms the hybrid/cosine
baseline.

% ===========================================================================
\section{Diversity-aware Ranking with MMR}

Maximal Marginal Relevance is implemented in \texttt{app.py}:
\begin{equation}
\text{MMR} = \lambda \cdot \text{relevance}
- (1-\lambda)\cdot \max_{\text{selected}} \text{similarity}.
\end{equation}
In the live recommendation flow, LTR first selects the top 100 candidates, then
MMR iteratively chooses 10 papers with $\lambda = 0.7$ to balance relevance and
redundancy. MMR is therefore part of the actual pipeline, but the proposal's
\emph{Naive Top-K vs MMR} ablation measured by Intra-list Diversity (ILD) is not
yet tabulated, so the diversity improvement from MMR is not quantified.

% ===========================================================================
\section{Explainable Recommendations and Interface}

The Streamlit UI in \texttt{app.py} supports selecting TF-IDF or SBERT,
choosing an existing or cold-start user, viewing reading history, generating
Top-10 recommendations, and displaying, per recommendation, the category, top
TF-IDF keywords, LTR score, semantic similarity, and recency score. It also
shows global LTR feature importance and visualises the user profile and
recommendations with t-SNE.

Relative to the proposal, most of the explanation information is present. The
main gap is per-recommendation feature contribution: the UI currently shows
only the model's global feature importance.

% ===========================================================================
\section{Evaluation Strategy and Results}

Evaluation is run on the 7{,}701-paper test split using 500 synthetic test
users, random seed 42, half-life 90, hybrid $\alpha=0.7$, and $K\in\{5,10,20\}$.
The metrics are Category Precision@K, NDCG@K, Hit Rate@K, Semantic Coherence,
Intra-list Diversity, and Silhouette Score.

\safefigure{eval/outputs/representation_metrics.png}{Representation comparison metrics.}{fig:repr-metrics}

\begin{table}[H]
\centering
\caption{Representation comparison across $K$.}\label{tab:repr-metrics}
\small
\begin{tabular}{lrrrrrr}
\toprule
\textbf{Method} & \textbf{K} & \textbf{Cat.\ Prec.} & \textbf{NDCG} & \textbf{Hit Rate} & \textbf{Sem.\ Coh.} & \textbf{ILD}\\
\midrule
TF-IDF & 5  & 0.8276 & 0.0000 & 0.000 & 0.2751 & 0.7934\\
TF-IDF & 10 & 0.8164 & 0.0006 & 0.002 & 0.2662 & 0.8106\\
TF-IDF & 20 & 0.8068 & 0.0011 & 0.004 & 0.2563 & 0.8299\\
SBERT pretrained & 5  & 0.8848 & 0.0000 & 0.000 & 0.7305 & 0.3830\\
SBERT pretrained & 10 & 0.8796 & 0.0000 & 0.000 & 0.7229 & 0.3960\\
SBERT pretrained & 20 & 0.8720 & 0.0000 & 0.000 & 0.7138 & 0.4095\\
SBERT fine-tuned & 5  & 1.0000 & 0.0000 & 0.000 & 0.7997 & 0.1696\\
SBERT fine-tuned & 10 & 1.0000 & 0.0000 & 0.000 & 0.7910 & 0.1826\\
SBERT fine-tuned & 20 & 1.0000 & 0.0000 & 0.000 & 0.7803 & 0.2008\\
\bottomrule
\end{tabular}
\end{table}

Fine-tuned SBERT achieves the highest Category Precision and Semantic
Coherence, whereas TF-IDF achieves the highest ILD but the lowest coherence.
This illustrates the relevance--diversity trade-off and motivates applying MMR
on top of fine-tuned SBERT. NDCG and Hit Rate are very low because the target
is a single specific paper while the system ranks over all 7{,}701 candidates:
many other papers may match the topic without being counted as the positive.
These two metrics are therefore not strong enough to draw conclusions about
target retrieval.

\safefigure{eval/outputs/tsne_compare_all.png}{t-SNE comparison of the three representations.}{fig:tsne}

\begin{table}[H]
\centering
\caption{Silhouette Score by representation.}\label{tab:silhouette}
\begin{tabular}{lr}
\toprule
\textbf{Method} & \textbf{Silhouette Score}\\
\midrule
TF-IDF & 0.0153\\
SBERT pretrained & 0.0467\\
SBERT fine-tuned & \textbf{0.3749}\\
\bottomrule
\end{tabular}
\end{table}

Fine-tuned SBERT forms clearly separated category clusters: its Silhouette
Score is roughly 8$\times$ that of pretrained SBERT and about 24$\times$ that of
TF-IDF. The t-SNE plots are illustrative only, since a 2-D projection can
distort distances; Silhouette Score and the ranking metrics are the primary
quantitative evidence.

% ===========================================================================
\section{Mapping to Proposal Objectives}

\begin{table}[H]
\centering
\caption{Status of each proposal objective.}\label{tab:status}
\small
\begin{tabular}{p{3.1cm}p{2.4cm}p{4.2cm}p{4.0cm}}
\toprule
\textbf{Item} & \textbf{Status} & \textbf{Evidence} & \textbf{Missing}\\
\midrule
Data preprocessing & Basic complete & Train/test CSV, preprocessor, recency score & Some proposal metadata missing\\
TF-IDF baseline & Complete & Train/test matrix, profile, evaluation & ---\\
SBERT pretrained & Complete & Train/test embeddings, evaluation & ---\\
Fine-tuned SBERT & Artifacts complete & Checkpoints, embeddings, t-SNE, metrics & Triplet pipeline not fully reproducible\\
Temporal user profile & Complete & 1{,}000 users, half-life 30/60/90 & Mean-vs-decay ablation missing\\
Cold-start & Complete & Code, UI, unit tests & No dedicated metric yet\\
Hybrid scoring & Complete & TF-IDF/SBERT results, three half-lives & Needs semantic/recency normalisation\\
Learning-to-Rank & Basic complete & Two models, UI integration & Two features and baseline benchmark missing\\
MMR & Integrated & Reranking in Streamlit & Naive Top-K vs MMR ablation missing\\
Explainability & Mostly complete & Keywords and score breakdown & Per-item LTR contribution missing\\
Streamlit demo & Present & \texttt{app.py} & Needs stable end-to-end testing\\
Evaluation framework & Complete for representation & 6 metrics, CSV/JSON/PNG & LTR, MMR, profile ablations not evaluated\\
\bottomrule
\end{tabular}
\end{table}

% ===========================================================================
\section{Testing}

\begin{table}[H]
\centering
\caption{Test results.}\label{tab:tests}
\begin{tabular}{lr}
\toprule
\textbf{Test group} & \textbf{Result}\\
\midrule
Core project tests & 15/15 pass\\
Evaluation tests & 10/10 pass\\
Total & \textbf{25/25 pass}\\
\bottomrule
\end{tabular}
\end{table}

The tests cover cold-start, user profile generation, user t-SNE, synthetic test
user creation, and the metric functions.

% ===========================================================================
\section{Limitations and Remaining Work}

To fully satisfy the expected outcomes in the proposal, four groups of results
are still missing:
\begin{enumerate}[leftmargin=1.6em]
  \item \textbf{Cosine baseline vs LTR:} run on the same candidate pool and user
  set; report NDCG, Hit Rate, and Semantic Coherence; produce an official
  feature-importance plot.
  \item \textbf{Naive Top-K vs MMR:} report ILD before and after MMR and measure
  the relevance drop to quantify the trade-off.
  \item \textbf{Mean profile vs temporal-weighted profile:} compare mean and
  half-life 30/60/90 using Category Precision, NDCG, and Semantic Coherence.
  \item \textbf{Reproducibility:} store the triplet file (or a runnable triplet
  generator with current paths), align the fine-tuning method between proposal
  and implementation, and save evaluation output for LTR and MMR.
\end{enumerate}

In addition, the hybrid score needs calibration, since recency currently
contributes only about 3\% with SBERT despite a configured weight of 30\%.

% ===========================================================================
\section{Conclusion}

The repository realises a fairly complete recommendation pipeline in the
direction of the proposal: preprocessing, three representations, a
temporal user profile, cold-start handling, hybrid scoring, LTR, MMR,
explainability, and a Streamlit UI all have implementations or stored
artifacts.

The strongest quantitative results lie in the representation stage:
fine-tuned SBERT reaches a Category Precision of $1.0000$, a Semantic Coherence
of $0.7910$ at $K=10$, and a Silhouette Score of $0.3749$, clearly above both
baselines. However, the full proposal has not yet been confirmed
experimentally: three important ablations --- LTR, MMR, and the temporal
profile --- still lack comparison tables. The fair status of the project is thus
that \emph{the pipeline is largely implemented, but the evaluation matrix in the
proposal is not yet complete}.

% ===========================================================================
\section{Key Artifacts}

\begin{table}[H]
\centering
\caption{Main project artifacts.}\label{tab:artifacts}
\begin{tabular}{ll}
\toprule
\textbf{Component} & \textbf{Path}\\
\midrule
Original proposal & \texttt{proposal.docx}\\
Streamlit application & \texttt{app.py}\\
Cold-start & \texttt{cold\_start.py}\\
User profile & \texttt{user\_profile\_generation.py}\\
LTR training & \texttt{train\_ltr.py}, \texttt{LTR.ipynb}\\
Fine-tuning & \texttt{finetune\_BERT/}\\
Evaluation & \texttt{eval/evaluate\_representations.py}\\
Synthetic test users & \texttt{eval/generate\_test\_users.py}\\
Metric source & \texttt{eval/outputs/representation\_metrics.json}\\
Metric chart & \texttt{eval/outputs/representation\_metrics.png}\\
t-SNE chart & \texttt{eval/outputs/tsne\_compare\_all.png}\\
\bottomrule
\end{tabular}
\end{table}

% ===========================================================================
\begin{thebibliography}{9}
\bibitem{kaggle} ``7700+ Latest ArXiv AI/ML Research Papers (2025--2026),''
Kaggle dataset.
\bibitem{sbert} N.~Reimers and I.~Gurevych, ``Sentence-BERT: Sentence
Embeddings using Siamese BERT-Networks,'' \emph{EMNLP}, 2019.
\bibitem{minilm} W.~Wang et al., ``MiniLM: Deep Self-Attention Distillation for
Task-Agnostic Compression of Pre-Trained Transformers,'' \emph{NeurIPS}, 2020.
\bibitem{mmr} J.~Carbonell and J.~Goldstein, ``The Use of MMR, Diversity-Based
Reranking for Reordering Documents and Producing Summaries,'' \emph{SIGIR},
1998.
\bibitem{gbm} J.~H. Friedman, ``Greedy Function Approximation: A Gradient
Boosting Machine,'' \emph{Annals of Statistics}, vol.~29, no.~5, 2001.
\bibitem{ndcg} K.~J\"arvelin and J.~Kek\"al\"ainen, ``Cumulated Gain-Based
Evaluation of IR Techniques,'' \emph{ACM TOIS}, vol.~20, no.~4, 2002.
\bibitem{sklearn} F.~Pedregosa et al., ``Scikit-learn: Machine Learning in
Python,'' \emph{JMLR}, vol.~12, 2011.
\bibitem{streamlit} Streamlit, ``Streamlit API Reference,''
\url{https://docs.streamlit.io/}.
\end{thebibliography}

\end{document}
